% AUTO-GENERATED by the update_record tool. DO NOT EDIT.
% verify_paper regenerates this file from record.json and the run outputs
% and fails on any difference, so manual edits always surface.
\noindent\textbf{Hypothesis 1.} SciGym 公開コード（コミット 88a7b93）の Controller をそのまま用い、Vercel AI Gateway の google/gemini-3.8-flash に最大 5 反復・temperature 1.0 で SciGym-small の 10 件（pilot）を解かせたとき、平均 Simulation Trajectory Error（STE）は論文 Table 1 の Gemini-2.5-Flash（0.4181）以下である。~\cite[s1.p1, s1.p3]{duan-2025-measuring}
\begin{enumerate}
\item[\textbf{Claim 1}] google/gemini-3.8-flash で SciGym-small の 10 件を最大 5 反復で解かせたときの平均 STE は、論文 Table 1 の Gemini-2.5-Flash の 0.4181 以下である。~\cite[s1.p1, s1.p3, s1.p4]{duan-2025-measuring}
  \emph{Rationale:} 論文の 1 世代前の Flash を現行の Flash が上回れば、公式 Controller を無改変で走らせた上で h1 の数値部分が成り立つ。
  \emph{Criterion:} \texttt{\detokenize{gemini-3.8-flash}}.\texttt{\detokenize{trajectory_smape}} $-$ 0.4181 $\leq 0$.
  \emph{Prediction:} $[-0.3, 0]$ (同モデルは 137 件・20 反復で 0.218 だった（auto-res2/scigym-frontier-models）。反復を 5 に減らしても 0.4181 は下回ると予測する).
  \emph{Observed:} -0.28363.
  \emph{Verdict:} supported.
\end{enumerate}
\noindent\emph{Assumed, so that the claims together imply Hypothesis 1:}
\begin{itemize}
\item 10 件の平均 STE が 137 件の平均を代表すること。本研究は record の宣言と実行時観測の照合（repository\_integration と observed.json）を確かめる統合確認であり、件数は小さい（c1）
\item Vercel AI Gateway 経由の gemini-3.8-flash が、同じモデルを直接呼んだ場合と同じ能力で応答すること（c1）
\end{itemize}
